\renewcommand*{\authorOfNextLines}{Helmut Quinz}
Im Hauptdokument finden Sie die Definition\\
\verb|\newif\ifCompileFinal\CompileFinalfalse|\\
Damit wird festgelegt, ob das Dokument als noch zu bearbeitende Version übersetzt werden soll oder als final abzugebend.
In ersterem Fall zeigt der Compiler durch schwarze Rechtecke welchen Stellen der Text über die festgelegte Breite hinaus gesetzt wurde.
Hier ist in der Textkomposition manuell nachzuarbeiten.\\
Zusätzlich zeigt so das Dokument im Anhang die vom Paket \texttt{fixme} erkannten Anmerkungen an.\\

Die Umstellung zwischen den Versionen erfolgt über die Angabe 
\begin{description}
	\item[\texttt{\textbackslash CompileFinalfalse}] für eine noch zu bearbeitende Version \bzw
	\item[\texttt{\textbackslash CompileFinaltrue}] für die finale Version 
\end{description}
\subsubsection{Hinweis}
Findet der Compiler mit \verb|\fxfatal| definierte Anmerkungen lässt sich das Dokument nicht in seiner finalen Version erstellen!